%% sections/02_background.tex

\section{배경}
\label{sec:background}

\subsection{투기적 디코딩과 드래프터}
\label{subsec:specdecode-bg}

투기적 디코딩은 드래프터가 제안한 $\gamma$개의 후보 토큰을
타깃 모델이 단일 forward pass로 검증하고,
거부 위치 이전까지를 수락한다~\cite{leviathan2023spec,chen2023spec}.
드래프터는 (a) 별도 소형 모델, (b) 추가 헤드(Medusa~\cite{cai2024medusa},
EAGLE~\cite{li2024eagle,li2025eagle3}), (c) 모델 없는 검색
(Prompt Lookup n-gram~\cite{saxena2024prompt},
SuffixDecoding~\cite{snowflake-arctic})으로 나뉜다.
본 연구는 주 드래프터로 SuffixDecoding(suffix tree 기반)을 사용하고,
naive n-gram 과 vanilla(투기 없음), 그리고 라우팅 베이스라인 llm-d 를 비교한다.
SuffixDecoding 은 입력과 \emph{이전 생성 출력}을 suffix tree 에 누적하여
프롬프트 외 반복 구조까지 드래프트에 활용한다.

\subsection{R/K 프레임워크와 수락률}
\label{subsec:rk-framework}

투기적 디코딩의 순이익은 \emph{수락률} $\alpha$ 와 \emph{step overhead} 의 함수다.
한 step 의 기대 수락 토큰 수를 $K_{\mathrm{eff}}$, draft+verify 의
step당 상대 비용을 overhead 라 하면, vanilla 대비 처리량 비는 근사적으로
\begin{equation}
\frac{T_{\mathrm{spec}}}{T_{\mathrm{vanilla}}}
  \;\approx\; \frac{1 + \alpha\,K_{\mathrm{eff}}}{1 + \mathrm{overhead}},
\label{eq:rk-ratio}
\end{equation}
이며, $\alpha\,K_{\mathrm{eff}} > \mathrm{overhead}$ 일 때만 가속이 된다.
$\alpha$ 는 타깃 모델이 드래프트 토큰을 수락하는 비율로,
드래프터 품질과 출력의 \emph{예측 가능성}에 의해 결정된다.
overhead 는 모델 규모와 아키텍처(특히 MoE 의 expert 활성 비용)에 비례한다.
식~\ref{eq:rk-ratio}은 부호와 방향에 대한 정성적 도구이며,
처리량의 정확한 크기는 본 연구에서 실측으로 보고한다.

\subsection{추론 특화 모델과 MoE}
\label{subsec:reasoning-moe-bg}

최근 fleet 에는 두 부류의 모델이 함께 배치된다.
\emph{추론 특화} 모델(DeepSeek-R1 계열~\cite{deepseekr1})은 답 이전에 긴
사고 사슬(chain-of-thought)을 생성하며, 같은 입력에도 다양한 추론 경로를 취한다.
\emph{Mixture-of-Experts}(MoE) 모델은 토큰마다 일부 expert 만 동적으로 활성화하여
파라미터 수 대비 활성 연산을 줄이되, 라우팅이 토큰 의존적이다.
두 특성 모두 출력 토큰 분포의 예측 가능성을 낮추는 방향으로 작용하며
(\S\ref{sec:mechanism}), 이는 드래프트 수락률 $\alpha$ 와 직결된다.
본 연구는 dense 7B--405B 와 MoE 671B(R1)를 함께 측정하여 이 효과를 분리한다.

\subsection{연속 배칭과 자원 이용률}
\label{subsec:cb-bg}

프로덕션 서빙은 연속 배칭(continuous batching~\cite{yu2022orca}) 하에서
다수 요청을 동시 처리한다. 따라서 처리량은 단일 요청 지연이 아니라
GPU step 점유와 host 측 step간 작업(스케줄링$\cdot$드래프팅$\cdot$샘플링$\cdot$detok)에
함께 좌우된다. 본 연구는 nvidia-smi 와 \texttt{/proc} 동기 샘플링(UtilSampler)으로
GPU$\cdot$CPU 이용률을 셀별로 수집하여 spec-decode 가 이 자원 균형을
어떻게 바꾸는지 측정한다(\S\ref{subsec:res-resource}).
